%
% Front matter
%

\begin{frontmatter}

\title{
  T2K 1.431$\times10^{20}$-POT Muon-Neutrino Disappearance Analysis
}

\vspace*{0.7cm}

\author[RAL]     {Costas Andreopoulos\corref{cor}}
\author[RAL,OXF] {Thomas Dealtry}
\author[ICL,RAL] {James Dobson}
\author[LANCS]   {Nick Grant}
\author[OXF]     {Martin Haigh}
\author[RAL]     {Jelena Ilic}
\author[RAL]     {Geoff Pearce}
\author[RAL,OXF] {and Alfons Weber}

\cortext[cor]{E-mail: costas.andreopoulos@stfc.ac.uk}

\address[RAL]{Particle Physics Dept., Rutherford Appleton Laboratory, STFC \\
              Harwell Science and Innovation Campus\\
              Oxfordshire OX11 0QX, UK}
\address[OXF]{Subdepartment of Particle Physics, University of Oxford \\
              Denys Wilkinson Building, Keble Road\\
              Oxfordshire OX1 3RH, UK}
\address[LANCS]{Physics Dept., Lancaster University \\
              Lancaster, LA1 4YB, UK}
\address[ICL]{Physics Dept., Imperial College London \\
              Blackett Lab, Prince Consort Road\\
              London SW7 2BW, UK}

%
% Abstract
%

\begin{abstract}

We report the results of a 2-flavour $\nu_{\mu}$-disappearance analysis on the combined Run1+2 dataset.
This combined dataset corresponds to an integrated JPARC neutrino beam
exposure of $1.431 \times 10^{20}$ POT.
The oscillation analysis was performed using both GENIE \cite{Genie:NIM} and NEUT \cite{Neut:ActaP}.
All of the inputs to the analysis are official 2010 inputs.

Using the ND280 measurement, GENIE predicts 101.72 $\pm$ 10.09 (stat) $^{+13.88}_{-13.49}$ (syst)
1 $\mu$-like ring events in SuperK without oscillations but only 31 were observed.
The statistical significance of the observed deficit is $\sim$4.2 standard deviations.
NEUT predicts 103.59 $\pm$ 10.18 (stat) $^{+13.77}_{-13.43}$ (syst) events in SuperK,
which implies that the statistical significance
of the observed deficit is $\sim$4.3 standard deviations.

The observed number of 1$\mu$-like ring events is consistent with 
the 2011 MINOS best-fit $\nu_{\mu}$-disappearance hypothesis
($sin^{2}2\theta_{23}$ = 1.0 and $|\Delta m^{2}_{23}|$ = 2.32$\times10^{-3}$ $eV^{2}/c^{4}$) \cite{Minos2011}.
Under this oscillation hypothesis,
GENIE predicts 28.39 $\pm$ 5.33 (stat) $^{+4.53}_{-4.46}$ (syst) 1 $\mu$-like ring events
in SuperK, while NEUT predicts 29.18 $\pm$ 5.40 (stat) $^{+4.31}_{-4.26}$ (syst) events. 

A 2-flavour $\nu_{\mu}$-disappearance fit to
the observed reconstructed energy spectrum of 1 $\mu$-like ring events of the combined Run1+2 dataset,
with $sin^{2}2\theta_{23}$ constrained to the physical region,
yields
$|\Delta m^{2}_{23}|$ = 2.60$\times$10$^{-3}$ $eV^{2}/c^{4}$ and
$sin^{2}2\theta_{23}$ = 0.97
as the best-fit values for GENIE ($\chi^{2}/ndf$ = 0.310).
For NEUT, the corresponding values are
$|\Delta m^{2}_{23}|$ = 2.65$\times$10$^{-3}$ $eV^{2}/c^{4}$ and
$sin^{2}2\theta_{23}$ = 0.98 ($\chi^{2}/ndf$ = 0.311).

\end{abstract}

%
% Keywords
%

%\begin{keyword}
%\end{keyword}

\end{frontmatter}

\clearpage

\tableofcontents

\clearpage

\listoffigures
\listoftables


